\chapter{Mixing optimization in HM model}\label{app:sc-g-opt}

This short appendix is meant to give a practical example on a topic we often referenced to in text: how optimisation of $g$, the mixing parameter defined in Sec.~\ref{subsec:reject-mix-iterate} and given operatively in  Eq.~\eqref{eq:mix} as
\[
	\mathbf{v}_{i+1} \equiv g \mathrm{A} \left( \mathbf{v}_{i} \right) + (1-g) \mathbf{v}_{i}
\]
can be relevant in order to ensure convergence of an iterative HF scheme as presented in Sec.~\ref{sec:sketch-ihf-algorithm}.

\section{Map analysis for $s$-wave superconductivity in HM}

As we discussed in Sec.~\ref{subsec:sc-hm}, conventional repulsive HM cannot exhibit BCS $s$-wave superconductivity within a MFT general scheme. This is proven at the end of Sec.~\ref{subsec:sc-hm}, where we end up withe the self-consistency equation
\begin{equation}\label{appeq:sc-self-consistency}
	w^{(s)} = -\frac{U}{U_\mathrm{c}} w^{(s)}
\end{equation}
only solved at $U>0$ by $w^{(s)}=0$. Now, our aim is the following: given any initializer $w^{(s)}_0$, our iterative HF algorithm should give
\[
	w^{(s)}_0
	\quad\stackrel{\mathrm{Alg}}{\to}\quad
	\cdots 
	\quad\stackrel{\mathrm{Alg}}{\to}\quad 
	\sim 0 
\]
The point is the following: is this behaviour possible at any model parametrization?

Consider the nonlinear map of Eq.~\eqref{appeq:sc-self-consistency} as ``seen by the algorithm'': at step $i$, we compute the left side by using the value of $w^{(s)}$ at step $i-1$ (also for computing $U_\mathrm{c}[\beta]$). Let us drop the superscript $(s)$ for the sake of readability. Seen this way, we are effectively moving along the numeric succession
\[
	w_{i+1} = -g\left(\frac{U}{U_{\mathrm{c}}[\beta]}\right)_i w_i + (1-g) w_i \\
\]
the subscript indicating iteration. In previous equation we made use of Eq.~\eqref{eq:mix} to generate the $i+1$-th step from step $i$. Let now:
\[
\begin{aligned}
	U_\mathrm{s} &\equiv \left(\frac{1}{g}-1\right) U_{\mathrm{c}} \\
	U_\mathrm{h} &\equiv \left(\frac{2}{g}-1\right) U_{\mathrm{c}}
\end{aligned}
\]
the subscript indicating, for a reason clear soon, a ``soft boundary'' (first) and an ``hard boundary'' (second). Let us derive some trivial facts true whenever $w_i>0$:
\begin{enumerate}
	\item It always holds $w_{i+1} \le w_i$. This is simple to see: if $w_{i+1}>w_i$,
	\[
	\begin{aligned}
		-g\left(\frac{U}{U_{\mathrm{c}}[\beta]}\right)_i w_i + (1-g) w_i &> w_i \\
		-g\left(\frac{U}{U_{\mathrm{c}}[\beta]}\right)_i w_i -g w_i &> 0
	\end{aligned}
	\]
	which is impossible because all involved terms are positive.
	%
	\item For $0 < U < U_\mathrm{s}$, we have $0 < w_{i+1} < w_i$. Indeed $w_{i+1}$ is greater than zero when:
	\[
	\begin{aligned}
		-g\left(\frac{U}{U_{\mathrm{c}}}\right)_i w_i + (1-g) w_i &> 0 \\
		\frac{1-g}{g} &> \frac{U}{U_{\mathrm{c}}} \\
		U_\mathrm{s} &> U 
	\end{aligned}
	\]
	%
	\item For $U_\mathrm{s} < U < U_\mathrm{h}$, we have $-w_i < w_{i+1} < 0$. From point above, we know $U>U_\mathrm{s}$ implies $w_{i+1} < 0$ by simply inverting inequalities. Moreover,
	\[
	\begin{aligned}
		-g\left(\frac{U}{U_{\mathrm{c}}}\right)_i w_i + (1-g) w_i &> -w_i \\
		\frac{2-g}{g} &> \frac{U}{U_{\mathrm{c}}} \\
		U_\mathrm{h} &> U 
	\end{aligned}	
	\]
	%
	\item For $U > U_\mathrm{h}$, we have $w_{i+1} < -w_i$. This is trivially implied by inverting inequalities in last point.
\end{enumerate}
If we know $U_\mathrm{c}$, for any given $g$ we can predict the algorithmic image of a given step.

Now, consider $U_\mathrm{c}$ as a function of the gap $\Delta$,
\[
	U_\mathrm{c}(\Delta) = \left[
		\frac{1}{2 L_x L_y} \sum_{\mathbf{k}\in\mathrm{BZ}} \frac{1}{\sqrt{\xi_\mathbf{k}^2 + \abs{\Delta}^2}} \tanh(
			\frac{\beta}{2} \sqrt{\xi_\mathbf{k}^2 + \abs{\Delta}^2}
		)
	\right]^{-1}
\]
\todo